\documentclass[10pt,a4paper]{amsart}
\usepackage[margin=19mm]{geometry}
\usepackage{amsmath,amssymb,mathtools}
\usepackage[scr=rsfs,cal=euler]{mathalfa}
\usepackage{xfrac}
\usepackage[unicode,hidelinks,bookmarks=false]{hyperref}
\newcommand{\ie}{i.e.\ }
\newcommand{\cf}{cf.\ }
\newcommand{\resp}{respectively\ }
\newcommand{\Subsection}{\subsection}
\providecommand{\nobreakdash}{\nobreak\hbox{-}}
\setlength{\emergencystretch}{2em}
\multlinegap0pt
\usepackage{bm}
\usepackage{bbm}
\usepackage{dsfont}
\usepackage{xspace}
\usepackage{cancel}
\usepackage{mathtools}
\usepackage{amsthm}
\makeatletter
\@ifundefined{theorem}{\newtheorem{theorem}{Theorem}}{}
\makeatother
\title[On a Weiss-type Almost Monotonicity Formula]{On a Weiss-type Almost Monotonicity Formula}
\author{Aelson Sobral}
\date{Source: arXiv:2604.12712v2 (CC BY 4.0)}
\begin{document}
\maketitle

\noindent\emph{Source and licence.} On a Weiss-type Almost Monotonicity Formula. Version arXiv:2604.12712v2 (CC BY 4.0), distributed under the Creative Commons Attribution 4.0 International licence, as recorded in the arXiv licence metadata for this exact version and shown on the abstract page https://arxiv.org/abs/2604.12712v2. Full rights notice: licenses/arxiv-2604.12712-CC-BY-4.0.txt.\par\smallskip

\noindent\textit{Editorial scope.} Counted unit: the Theorem whose source label is \texttt{thm : almost mon formula} in the article (source0.tex lines 1162--1179, source label \texttt{thm : almost mon formula}), delivered together with the article's own complete original proof of it (source0.tex lines 1181--1318, one \texttt{proof} environment of 138 lines). No mathematics is written, repaired or re-derived: statement and proof are the article's own text line for line. Required context retained uncounted: func:AC (lines 1156--1158). Every definition, display and label of those lines is retained unchanged. Omitted from this package and not redistributed: 1--1155, 1159--1161, 1180--1180, 1319--1478. Nothing inside a retained excerpt is omitted or altered: every retained body is delivered exactly as the article's lines are written, so no omission receipt of any kind accompanies this package. Citations: the retained excerpts carry no citation command at all, so no bibliography entry of the article is transcribed and no reference is redistributed. Reference adaptation: none was needed and none was made; every reference inside the delivered unit resolves inside it, so no unresolved reference or citation is produced. Printed slips kept literal: no printed word, symbol, hypothesis or number of the counted statement or of its proof was altered. Layout and class: the article's own class is \texttt{amsart} with packages including amsmath, amssymb, amsthm, mathtools, mathrsfs, hyperref, cleveref, xcolor, tikz, verbatim, setspace, ulem; the wrapper is the lane's pinned \texttt{amsart} editorial wrapper at 10pt on A4, which restates the theorem-like environments the retained text needs and adds the article's own macro definitions that the retained text needs (none), with extra wrapper packages (bm, bbm, dsfont, xspace, cancel, mathtools). The delivered numbering is the unit's own. Assets: no retained span contains a figure environment, a graphics-inclusion command, a PDF-page inclusion, a table or a TikZ picture; no image file is copied into this package, the package declares no asset dependency and no image file is redistributed with it. Licence: version arXiv:2604.12712v2 of this article is distributed under the Creative Commons Attribution 4.0 International licence, as recorded in the arXiv licence metadata for this exact version and shown on the abstract page https://arxiv.org/abs/2604.12712v2. A full rights notice is delivered at licenses/arxiv-2604.12712-CC-BY-4.0.txt. Characters: every retained excerpt is pure ASCII, so the delivered engine, XeLaTeX with its default Latin Modern text font, reports no missing character.
\begin{equation}\label{func:AC}
    \mathcal J(u) \coloneqq \int \left( \frac{1}{2} \, \langle A(x)\nabla u, \nabla u \rangle + Q(x)\mathcal{X}_{\{u>0\}}\right)\,dx.
\end{equation}

\begin{theorem}\label{thm : almost mon formula}
Let $u \geq 0$ be an almost minimizer to the functional \eqref{func:AC} in $B_2$. Assume further that
\[
	\omega_A(\cdot) + \omega_Q(\cdot) + \varrho(\cdot) \in C^{0,dini} \qquad \text{and} \qquad A(x_0) = I_d,
\]
for $x_0 \in F(u)\cap B_1$, and that
\[
	\int_{B_r(x_0)} |Du|^2\,dx \lesssim r^d, \quad \text{for} \quad r<1.
\]  
Then,
\[
    \mathcal{A}_{u,x_0}(r) \coloneqq r^{-d}\mathcal{J}_{x_0}(u,B_r(x_0)) - \frac{1}{2}r^{-d - 1}\int_{\partial B_r(x_0)}u^2\, d\mathcal{H}^{d-1},
\]
is almost monotone, in the sense that there exists a nonnegative real function $g \in L^1(0,1)$ such that
$$
	\frac{d}{dr} \mathcal{A}_{u,x_0}(r) \geq -g(r).
$$
\end{theorem}

\begin{proof}
Fix $x_0\in F(u)$ and assume, without loss of generality, that $x_0=0$ and $B_r\Subset B_1$ for every $r\in(0,1)$.
Set
\[
    A_0 \coloneqq A(0) = I_d, \qquad Q_0\coloneqq Q(0)>0,
\]
and define the \emph{frozen} functional at $0$ by
\[
    \mathcal J_{0}(v,E)\coloneqq \int_{E}\left(\frac12|\nabla v|^2+Q_0\,\chi_{\{v>0\}}\right)\,dx, \qquad E\subset \mathbb R^d.
\]
Then the Weiss-type quantity becomes
\[
    \mathcal A_{u,0}(r)=r^{-d}\mathcal J_0(u,B_r)-\frac12\,r^{-d-1}\int_{\partial B_r}u^2\,d\mathcal H^{d-1}.
\]
We calculate $\frac{d}{dr} A_{u,0}(r)$. Using the co-area formula, we have
\[
    \frac{d}{dr}\,\mathcal J_0(u,B_r)=\mathcal J_0(u,\partial B_r)
    \coloneqq \int_{\partial B_r}\left(\frac12|\nabla u|^2+Q_0\,\chi_{\{u>0\}}\right)\,d\mathcal H^{d-1},
\]
and so a direct computation yields
\begin{align*}
    \frac{d}{dr}\mathcal A_{u,0}(r)
    &= -d\,r^{-d-1}\mathcal J_0(u,B_r)+r^{-d}\mathcal J_0(u,\partial B_r) \\
    &\quad +\frac12(d+1)\,r^{-d-2}\int_{\partial B_r}u^2\,d\mathcal H^{d-1}
    - r^{-d-1}\int_{\partial B_r}u\,\partial_\nu u\,d\mathcal H^{d-1}.
\end{align*}
On $\partial B_r$, decompose $\nabla u = (\partial_\nu u)\,\nu + D_\tau u$, so that $|\nabla u|^2=|\partial_\nu u|^2+|D_\tau u|^2$.
By elementary algebra,
\[
    \frac12|\partial_\nu u|^2-\frac{u\,\partial_\nu u}{r}+\frac12\frac{u^2}{r^2}
    =\frac12\Big(\partial_\nu u-\frac{u}{r}\Big)^2.
\]
Collecting terms, we obtain
\begin{equation}\label{eq:Weiss-derivative}
    \frac{d}{dr}\mathcal A_{u,0}(r) = \mathcal B(u)-d\,r^{-d-1}\mathcal J_0(u,B_r)+\mathcal H(u),
\end{equation}
where
\begin{align*}
\mathcal B(u)
&\coloneqq r^{-d}\int_{\partial B_r}\Big(\tfrac12|D_\tau u|^2+Q_0\,\chi_{\{u>0\}}\Big)\,d\mathcal H^{d-1}
+\frac12\,r^{-d-1}\int_{\partial B_r}u^2\,d\mathcal H^{d-1},\\
\mathcal H(u)
&\coloneqq \frac12\,r^{-d-1}\int_{\partial B_r}\Big(\partial_\nu u-\frac{u}{r}\Big)^2\,d\mathcal H^{d-1}\ge 0.
\end{align*}
Let $w$ be the $1$-homogeneous extension of $u|_{\partial B_r}$ to $B_r$, namely
\[
    w(x)\coloneqq \frac{|x|}{r}\,u\!\left(r\frac{x}{|x|}\right)\quad \text{for} \quad x \not = 0, \quad \text{and}\quad w(0)\coloneqq 0.
\]
Then, since $w$ is $1$-homogeneous, $\mathcal A_{w,0}(s)$ is constant for $0<s\le r$; in particular,
\[
    0 = \frac{d}{ds}\mathcal A_{w,0}(s)\Big|_{s=r}.
\]
Applying \eqref{eq:Weiss-derivative} to $w$ and using that $\mathcal H(w) = 0$, we get
\[
    0 = \mathcal B(w)-d\,r^{-d-1}\mathcal J_0(w,B_r).
\]
Since $w=u$ on $\partial B_r$, we also have $\mathcal B(u)=\mathcal B(w)$, hence
\begin{equation}\label{eq:B-u-frozen-w}
\mathcal B(u)=d\,r^{-d-1}\mathcal J_0(w,B_r).
\end{equation}
Now we compare with the original functional. By the almost minimality of $u$, for every competitor $v$ with $v-u\in H^1_0(B_r)$,
\[
	\mathcal J(u,B_r)\le (1+\varrho(r))\,\mathcal J(v,B_r).
\]
Taking $v=w$, and using $\mathcal J\ge 0$, we obtain
\begin{equation}\label{eq:almost-min-lower}
	\mathcal J(w,B_r)\ge (1-\varrho(r))\,\mathcal J(u,B_r),
\end{equation}
since $(1+\varrho)^{-1}\ge 1-\varrho$ for $\varrho\ge 0$. Next, we compare 
\[
\begin{aligned}
	\mathcal J_0(v,B_r)-\mathcal J(v,B_r)
	&=\frac12\int_{B_r}\big\langle (I_d-A(x))\nabla v,\,\nabla v\big\rangle\,dx \\
	&\quad +\int_{B_r}\big(Q_0-Q(x)\big)\,\chi_{\{v>0\}}\,dx.
\end{aligned}
\]
Combining this with \eqref{eq:B-u-frozen-w} and \eqref{eq:almost-min-lower} gives
\begin{align*}
	\mathcal B(u)
	&= d\,r^{-d-1}\mathcal J_0(w,B_r) \\
	&= d\,r^{-d-1}\mathcal J(w,B_r)
	+d\,r^{-d-1}\big(\mathcal J_0(w,B_r)-\mathcal J(w,B_r)\big)\\
	&\ge d\,r^{-d-1}(1-\varrho(r))\,\mathcal J(u,B_r)
	+d\,r^{-d-1}\big(\mathcal J_0(w,B_r)-\mathcal J(w,B_r)\big).
\end{align*}
Insert this lower bound for $\mathcal B(u)$ into \eqref{eq:Weiss-derivative} to obtain
\begin{align}
\frac{d}{dr}\mathcal A_{u,0}(r)
&\ge \mathcal H(u)
-d\,r^{-d-1}\varrho(r)\,\mathcal J(u,B_r)
+d\,r^{-d-1}\big(\mathcal J(u,B_r)-\mathcal J_0(u,B_r)\big)\notag\\
&\quad +d\,r^{-d-1}\big(\mathcal J_0(w,B_r)-\mathcal J(w,B_r)\big). \label{eq:key-ineq}
\end{align}
Finally, we estimate the error terms. For $x\in B_r$, by continuity of $A$ and $Q$,
\[
\|A(x)-I_d\|\le \omega_A(r),
\qquad
|Q(x)-Q_0|\le \omega_Q(r).
\]
Hence, for any $v$,
\[
\big|\mathcal J(v,B_r)-\mathcal J_0(v,B_r)\big|
\le \frac12\,\omega_A(r)\int_{B_r}|\nabla v|^2\,dx + \omega_Q(r)\,|B_r|.
\]
Applying this to $v=u$ and $v=w$ and using \eqref{eq:key-ineq} yields
\begin{align*}
\frac{d}{dr}\mathcal A_{u,0}(r)
&\ge \mathcal H(u)
-d\,r^{-d-1}\varrho(r)\,\mathcal J(u,B_r)\\
& - C\,r^{-d-1}\omega_A(r)\!\int_{B_r}\big(|\nabla u|^2+|\nabla w|^2\big)\,dx
- C\,\frac{\omega_Q(r)}{r},
\end{align*}
for a dimensional constant $C$.

Finally, by our assumptions, we have
\[
\int_{B_r}|\nabla u|^2\,dx + \int_{B_r}Q(x)\chi_{\{u>0\}}\,dx \le C_0\,r^{d},
\qquad
\int_{B_r}|\nabla w|^2\,dx \le C_0\,r^{d}.
\]
Therefore,
\[
r^{-d-1}\varrho(r)\mathcal J(u,B_r)\le C\,\frac{\varrho(r)}{r},
\quad
r^{-d-1}\omega_A(r)\!\int_{B_r}\big(|\nabla u|^2+|\nabla w|^2\big)\,dx \le C\,\frac{\omega_A(r)}{r}.
\]
We conclude that
\[
\frac{d}{dr}\mathcal A_{u,0}(r)
\ge \mathcal H(u) - C\,\frac{\omega_A(r)+\omega_Q(r)+\varrho(r)}{r}
\ge -\,g(r),
\]
where
\[
g(r)\coloneqq C\,\frac{\omega_A(r)+\omega_Q(r)+\varrho(r)}{r}\ge 0.
\]
Since $\omega_A+\omega_Q+\varrho\in C^{0,\mathrm{dini}}([0,1])$, we have $\int_0^1 \frac{\omega_A(t)+\omega_Q(t)+\varrho(t)}{t}\,dt<\infty$, hence $g\in L^1(0,1)$. This proves the almost monotonicity of $\mathcal A_{u,0}(r)$.
\end{proof}

\end{document}
